% Local working draft. All files in this directory are ignored by Git.
% Reviewer wording is preserved from reviewer_comments.txt; only formatting changes.
% Replace placeholders with responses and verified manuscript references.
% Compile from this directory with: pdflatex response_to_reviewers.tex
\documentclass[11pt,a4paper]{article}
\usepackage[T1]{fontenc}
\usepackage[utf8]{inputenc}
\usepackage[margin=25mm]{geometry}
\usepackage{xcolor}
\usepackage{graphicx}
\usepackage[hidelinks]{hyperref}

\definecolor{reviewerblue}{RGB}{0,65,125}
\newenvironment{reviewercomment}
  {\begin{quote}\small\color{reviewerblue}}
  {\end{quote}}
\newcommand{\placeholder}[1]{[#1]}
\setlength{\parindent}{0pt}
\setlength{\parskip}{0.6em}
\setlength{\emergencystretch}{2em}

\title{Response to the reviewers\\[0.5em]
\large Fast simulation of nonlinear deep resistive networks\\
for energy-based computation}
\author{}
\date{Working draft}

\begin{document}
\maketitle

\section*{Opening response}
\placeholder{Add an opening response and a summary of the completed revisions.}

\section*{Reviewer 1}
\subsection*{General remarks}
\begin{reviewercomment}
This manuscript presents a coordinate-descent (CD) solver for nonlinear
resistive networks with monotone nonlinear one-port elements. The central
thesis is that the primitive of a monotone current-voltage characteristic
is convex, as a consequence the network equilibrium can be formulated as
minimization of an electrical pseudo-energy. For the particular devices
they're studying the authors are able to reduce the solve step to a
coordinate update, in turn reduced to a scalar KCL solution. There is a
really nice trick where they apply lambert's W function and are able to
apply this for multiple types of nonlinearities. The reported speedups
over a commercial SPICE solver (ELDO) are very impressive!

I do feel that the emperical evidence backing the paper's claims are not
as strong as the mathematical content. As an example, the broader claim
of employing PWL to simulate a realistic device and demonstrate practical
scaling is not followed through on. The comparison against a
general-purpose SPICE simulator would have benefited from greater
analysis/isolation of confounders e.g., can we quantify how much of the
speed up is due to coordinate descent and how much (if any) could be due
to the use of a specialized equilibrium solver.

Detailed comments below:
\end{reviewercomment}

\noindent\textbf{Response:}
\placeholder{To be drafted.}

\noindent\textbf{Changes to the manuscript:}
\placeholder{Describe the revision and give its section, page, line, or figure reference, as applicable.}

\subsection*{Comment R1.1}
\begin{reviewercomment}
1. The authors indicate their work is applicable to ``realistic devices''
through the PWL formulation.

The authors indicate that measured tabulated data and their proposed PWL
implementation as being the primary means of leveraging their proposed
approach. However, the strongest PWL-experiment used weas not based on
measured devices or representative of a real device, instead it is
generated simply by sampling the double-Shockley curve at 200 points. But
the authors already offered an analytical solution to this, so this is
not strong enough evidence of the capability the authors are claiming.

At least one real device-derived test would considerably strengthen the
paper, even something derived through industry/foundry compact model
would be significantly more convincing here.
\end{reviewercomment}

\noindent\textbf{Response:}
\placeholder{To be drafted.}

\noindent\textbf{Changes to the manuscript:}
\placeholder{Describe the revision and give its section, page, line, or figure reference, as applicable.}

\subsection*{Comment R1.2}
\begin{reviewercomment}
2. The double-Shockley approximation is not justified well in text.

Equation (19) is an antiparallel diode equation, but the authors then
replace it with a significantly easlier approximation. This approximation
introduces its own error and that error is not quantified/studied. At the
moment, it is not quite clear to me whether the baseline SPICE simulations
use the full Eq. (19) or the approximated Eq. (21) one as their core solved
system? If SPICE instead uses Eq. (19), this should be stated explicitly.
In particular, it is extremely important that the Eq. (19)-to-Eq. (21)
approximation errors be quantified.
\end{reviewercomment}

\noindent\textbf{Response:}
\placeholder{To be drafted.}

\noindent\textbf{Changes to the manuscript:}
\placeholder{Describe the revision and give its section, page, line, or figure reference, as applicable.}

\subsection*{Comment R1.3}
\begin{reviewercomment}
3. What is the algorithmic benefit of coordinate descent?

The authors primary result is a speedup v.s., a general-purpose commercial
SPICE solver. However, the root cause of this speedup is not established
in the paper. As an example, a reasonable baseline might be to solve the
KCLs with a purpose-built Newton, damped-Newton, or related sparse solver.
After all the authors already use Newton-Raphson for the PWL updates and
to solve their Lambert-w formulation.
\end{reviewercomment}

\noindent\textbf{Response:}
\placeholder{To be drafted.}

\noindent\textbf{Changes to the manuscript:}
\placeholder{Describe the revision and give its section, page, line, or figure reference, as applicable.}

\subsection*{Comment R1.4}
\begin{reviewercomment}
4. Depth scaling is not shown

The Digits experiment only goes up to three hidden layers. The larger
MNIST demonstration is a 1568x100x20 network with limited hidden layers.
More crucially, the manuscript reports that netlists could not be
generated for the larger network. However, this is surprising because
netlist genration ought to be trivial here. More information on what is
occurring here would help disentangle convergence issues from simply
netlist generation? The paper motivates their work and indicates that
deeper networks require more iterations to reach a acceptable accuracy.
However, this work does not actually demonstrate results using even
9-layer "tiny" models?
\end{reviewercomment}

\input{r1_4_reply.tex}

\subsection*{Comment R1.5}
\begin{reviewercomment}
5. The reported SPICE netlist-generation failures require explanation.

The manuscript states that SPICE netlists could not be generated for the
two-hidden-layer width-1024 networks. This is surprising! Please identify
exactly which stage failed.
\end{reviewercomment}

\noindent\textbf{Response:}
\placeholder{To be drafted.}

\noindent\textbf{Changes to the manuscript:}
\placeholder{Describe the revision and give its section, page, line, or figure reference, as applicable.}

\subsection*{Comment R1.6}
\begin{reviewercomment}
6. The implemented Lambert-W update is not as \emph{straightforward} as the
text seems to imply.

The supplementary material explains that Lambert-W arguments are clipped
for numerical stability and the resulting estimate is refined with Newton
iterations. More importantly, the Newton residual itself appears to
evaluate a clipped exponential.
\end{reviewercomment}

\noindent\textbf{Response:}
\placeholder{To be drafted.}

\noindent\textbf{Changes to the manuscript:}
\placeholder{Describe the revision and give its section, page, line, or figure reference, as applicable.}

\subsection*{Comment R1.7}
\begin{reviewercomment}
7. Error characterization should better expose the difficult cases.

The main text emphasizes the P90 relative L1 error. The supplementary
material does provide P50, P90, and P99 statistics, and most configurations
show excellent agreement. However, some cases exhibit a noticeably larger
tail; for example, one three-hidden-layer single-Shockley configuration.

Since the SPICE-levels of accuracy is a primary contribution of this
paper, it would be helpful to capture tail-error characterization into
the main text.
\end{reviewercomment}

\noindent\textbf{Response:}
\placeholder{To be drafted.}

\noindent\textbf{Changes to the manuscript:}
\placeholder{Describe the revision and give its section, page, line, or figure reference, as applicable.}

\subsection*{Comment R1.8}
\begin{reviewercomment}
8. The mathematical presentation can be made considerably more accessible.

The mathematical exposition could be made much more straightfoward.
Figure 1 should appear before or alongside the general set notation and
should annotate the free/clamped nodes, conductances, nonlinear shunts,
and current directions used in the equations.

It would also help to motivate early on why certain approximations are
being made, before making them. Specific example would be to signpost
that a nondecreasing current law is being studied due to the
energy-formulation which will subsequently be discussed.
\end{reviewercomment}

\noindent\textbf{Response:}
\placeholder{To be drafted.}

\noindent\textbf{Changes to the manuscript:}
\placeholder{Describe the revision and give its section, page, line, or figure reference, as applicable.}

\clearpage
\section*{Reviewer 2}
\subsection*{General remarks}
\begin{reviewercomment}
\textbf{Summary}

This paper proposes a fast solver for nonlinear deep neural network (DNN)
models that predicts computing-array characteristics and integrates the
solver into the PyTorch framework. The approach aims to enable efficient
array-level modeling for hardware-aware DNN evaluation.

\textbf{Strengths}

\begin{itemize}

\item Detailed modeling: The paper provides a detailed and comprehensive
treatment of the target problem and the considered device/array behavior.

\item PyTorch integration: Integration with PyTorch is useful for
incorporating array-level effects into DNN evaluation and design
exploration.

\item Fast solver: The focus on computational efficiency is valuable for
reducing the cost of detailed array-level simulation.

\end{itemize}
\end{reviewercomment}

\noindent\textbf{Response:}
\placeholder{To be drafted.}

\noindent\textbf{Changes to the manuscript:}
\placeholder{Describe the revision and give its section, page, line, or figure reference, as applicable.}

\subsection*{Weaknesses}

\subsection*{Comment R2.1}
\begin{reviewercomment}
1. Limited scope: The approach appears narrowly tailored to a specific
type of device. The authors should discuss how readily the methodology
can be extended to other device technologies and array architectures.
\end{reviewercomment}

\noindent\textbf{Response:}
\placeholder{To be drafted.}

\noindent\textbf{Changes to the manuscript:}
\placeholder{Describe the revision and give its section, page, line, or figure reference, as applicable.}

\subsection*{Comment R2.2}
\begin{reviewercomment}
2. Insufficient array-level details: More information is needed on the
considered array architecture, including array size, device/resistance
values, weight storage and mapping, input/output representation, and how
computation and readout are performed. These details are important for
assessing the practical relevance of the proposed modeling approach.
\end{reviewercomment}

\noindent\textbf{Response:}
\placeholder{To be drafted.}

\noindent\textbf{Changes to the manuscript:}
\placeholder{Describe the revision and give its section, page, line, or figure reference, as applicable.}

\subsection*{Comment R2.3}
\begin{reviewercomment}
3. Limited workload evaluation: The evaluation appears to focus primarily
on MNIST. The authors should demonstrate the scalability of the approach
using more complex DNN workloads and report how runtime and accuracy
scale with array and model size.
\end{reviewercomment}

\noindent\textbf{Response:}
\placeholder{To be drafted.}

\noindent\textbf{Changes to the manuscript:}
\placeholder{Describe the revision and give its section, page, line, or figure reference, as applicable.}

\subsection*{Comment R2.4}
\begin{reviewercomment}
4. Incomplete discussion of prior work: The paper should discuss relevant
prior work, including GENIEx (Chakraborty et al, DAC 2020) and TWINN
(Malhotra et al, TCAS\_I, 2025), as well as other device/crossbar modeling
approaches. The authors should clearly distinguish the proposed approach
from these existing methods in terms of modeling methodology, accuracy,
computational cost, and scope.
\end{reviewercomment}

\noindent\textbf{Response:}
\placeholder{To be drafted.}

\noindent\textbf{Changes to the manuscript:}
\placeholder{Describe the revision and give its section, page, line, or figure reference, as applicable.}

\subsection*{Comment R2.5}
\begin{reviewercomment}
5. Limited quantitative comparison: A quantitative comparison with
relevant prior approaches would help establish the advantages and
novelty of the proposed solver.
\end{reviewercomment}

\noindent\textbf{Response:}
\placeholder{To be drafted.}

\noindent\textbf{Changes to the manuscript:}
\placeholder{Describe the revision and give its section, page, line, or figure reference, as applicable.}

\subsection*{Overall assessment}
\begin{reviewercomment}
The paper presents a useful and detailed approach for fast nonlinear
array modeling, with the PyTorch integration being particularly useful
for DNN evaluation. However, the current scope and evaluation appear
limited. Additional array-level details, evaluation on more complex
workloads, and a stronger comparison with prior work would strengthen
the paper and better establish its generality and practical applicability.
\end{reviewercomment}

\noindent\textbf{Response:}
\placeholder{To be drafted.}

\noindent\textbf{Changes to the manuscript:}
\placeholder{Describe the revision and give its section, page, line, or figure reference, as applicable.}

\end{document}
